\documentclass[11pt,a4paper]{article}
\usepackage[margin=24mm,headheight=15pt]{geometry}
\usepackage{fontspec}
\setmainfont{Times New Roman}
\setsansfont{Avenir Next}
\setmonofont{Menlo}
\usepackage{microtype}
\usepackage{newtxmath}
\usepackage{mathtools}
\usepackage{booktabs,longtable,array,tabularx}
\usepackage{pdflscape}
\usepackage[table]{xcolor}
\usepackage{enumitem}
\usepackage{fancyhdr}
\usepackage{titlesec}
\usepackage{xurl}
\usepackage{hyperref}
\usepackage{bookmark}
\usepackage{fancyvrb}
\usepackage[most]{tcolorbox}

\definecolor{Navy}{HTML}{18324A}
\definecolor{Accent}{HTML}{C56A3A}
\definecolor{LinkBlue}{HTML}{2D6A8A}
\definecolor{SoftBlue}{HTML}{EAF1F5}
\definecolor{TableHead}{HTML}{DCE8EF}
\definecolor{BodyGray}{HTML}{26343D}

\hypersetup{
  colorlinks=true,
  linkcolor=LinkBlue,
  urlcolor=LinkBlue,
  citecolor=LinkBlue,
  pdftitle={Proof audit of the graph-family notes},
  pdfauthor={Final Thesis Master Project}
}
\setlength{\parindent}{0pt}
\setlength{\parskip}{0.65em}
\setlength{\emergencystretch}{3em}
\setlist{leftmargin=1.6em,itemsep=0.3em,topsep=0.25em}
\setcounter{tocdepth}{2}
\renewcommand{\arraystretch}{1.15}
\allowdisplaybreaks

\titleformat{\section}
  {\Large\bfseries\sffamily\color{Navy}}
  {\thesection}{0.65em}{}
  [\vspace{-0.25em}\color{Accent}\titlerule]
\titleformat{\subsection}
  {\large\bfseries\sffamily\color{Navy}}
  {\thesubsection}{0.6em}{}
\titleformat{\subsubsection}
  {\normalsize\bfseries\sffamily\color{LinkBlue}}
  {\thesubsubsection}{0.55em}{}

\newtcolorbox{ResearchQuote}{
  enhanced,
  breakable,
  colback=SoftBlue,
  colframe=LinkBlue,
  boxrule=0pt,
  leftrule=2.2pt,
  arc=0pt,
  left=8pt,right=8pt,top=5pt,bottom=5pt
}
\newenvironment{ResearchContents}
  {\begin{tcolorbox}[breakable,colback=SoftBlue,colframe=SoftBlue,arc=2pt]}
  {\end{tcolorbox}}

\pagestyle{fancy}
\fancyhf{}
\fancyhead[L]{\small\sffamily\color{Navy} Asymmetric stochastic TSP}
\fancyhead[R]{\small\sffamily\color{Navy}\detokenize{audit.md}}
\fancyfoot[C]{\small\sffamily\color{BodyGray}\thepage}
\renewcommand{\headrulewidth}{0.3pt}
\renewcommand{\footrulewidth}{0pt}

\begin{document}
\color{BodyGray}
\begin{titlepage}
  \thispagestyle{empty}
  \vspace*{22mm}
  {\sffamily\small\bfseries\color{Accent}\MakeUppercase{Research note}}\par
  \vspace{8mm}
  {\sffamily\fontsize{21}{25}\selectfont\bfseries\color{Navy} Proof audit of the graph-family notes}\par
  \vspace{6mm}
  {\color{Accent}\rule{42mm}{1.4pt}}\par
  \vspace{11mm}
  {\Large\color{BodyGray} Asymmetric stochastic TSP\\
  clairvoyance-gap investigation}\par
  \vfill
  \begin{tcolorbox}[
    colback=SoftBlue,colframe=SoftBlue,arc=2pt,
    left=10pt,right=10pt,top=8pt,bottom=8pt]
    \sffamily
    \textbf{Source}\quad \detokenize{audit.md}\\[3pt]
    \textbf{Prepared}\quad 25 July 2026\\[3pt]
    \textbf{Format}\quad Typeset mathematical PDF
  \end{tcolorbox}
  \vspace{12mm}
\end{titlepage}

\begin{ResearchContents}
\tableofcontents
\end{ResearchContents}
\clearpage

\phantomsection
\section*{Scope and audit standard}
\addcontentsline{toc}{section}{Scope and audit standard}

This audit checks the completed notes against the stochastic-TSP model in \nolinkurl{problem-brief.md}.  In particular:


\begin{itemize}

\item a policy moves only when it calls an active client (apart from the final return to the depot);

\item an inactive call causes no movement;

\item a client passed through by an expanded shortest path is neither called nor served;

\item every client is called exactly once by an adaptive policy;

\item shortest-path walks may use inactive and uncalled client vertices as transit;

\item activations are mutually independent; and

\item local service bounds must survive arbitrary interleaving.

\end{itemize}

“Sound” below means that the stated result is justified at the level claimed. “Needs repair” means that the conclusion is likely correct but the written proof omits a model-specific step.  “Overstated” means that the displayed calculation is valid only for an extra conditional template or a particular candidate routing ledger, not for the stochastic-TSP optimum asserted or suggested by the surrounding text.


\phantomsection
\section*{Executive findings}
\addcontentsline{toc}{section}{Executive findings}

The following central results survive the audit:


\begin{enumerate}

\item the additive-replication inequality in \nolinkurl{composition-principles.md};

\item the exact \(8/7\) Cayley gadget in \nolinkurl{cayley-digraphs.md};

\item the exact ratio-one butterfly output-terminal construction in \nolinkurl{switching-networks.md};

\item the exact ratio-one projective-plane construction in \nolinkurl{projective-planes.md};

\item the potential-reweighting obstruction, the tree-tour bounds, and the conditional Hamiltonian ratio-one theorem in \nolinkurl{generalized-polygons.md};

\item the independent-codeword obstruction, Tanner potential obstruction, and conditional regular-design theorem in \nolinkurl{codes-designs-ltc.md};

\item the abstract product-choice permutation propositions in \nolinkurl{algebraic-lifts.md}; and

\item the Ramanujan preferred-orientation divergence calculation, when explicitly restricted to an augmentation which retains every preferred edge trace; and

\item the four-star layered-poset construction, including its interleaved Bellman product lemma, positive directed-metric embedding, asymptotic liminf \(3/2\) lower bound, and explicit finite \(>4/3\) certificate.

\end{enumerate}

There are two substantive qualifications.


\begin{itemize}

\item The algebraic-lift permutation theorem is correct, but the subsequent \(\tfrac74|\Gamma|w\) “posterior accounting” and \(8/7\) ratio are \textbf{not} bounds on the completed stochastic-TSP instance.  The displayed layer does not identify input and output ports, does not force every inactive preferred trace into an optimum tour, and supplies no adaptive upper bound of \(2|\Gamma|w\).  The only rigorous conclusion is that a particular all-preferred-traces circulation would have linear defects under additional port assumptions.

\item The high-girth scaffold additivity proposition is credible and repairable, but its adaptive upper-bound proof must use the pending-walk/policy-gluing argument from Chapter 4.  A policy cannot literally follow a scaffold walk, voluntarily return from a local module, and then enter the next one.

\end{itemize}

The four-star layered-poset note now proves a ratio above \(4/3\).  Earlier negative statements in this audit refer only to the original eight off-the-shelf family templates and not to this switching/design hybrid.


\phantomsection
\section*{0. Four-star layered poset}
\addcontentsline{toc}{section}{0. Four-star layered poset}

The complete proof and a \(120002\)-client witness are in \nolinkurl{four-star-layered-poset-construction.md}; an independently derived finite certificate is in \nolinkurl{four-star-poset-finite-certificate.md}.


\begin{itemize}

\item \textbf{Product law:} all \(4L\) selector clients have independent Bernoulli-\(q\) activations; the two gates per layer are permanent.

\item \textbf{Posterior:} every set of at most two row/column stars has complementary \(K_{2,2}\) representatives, and \(\mathbb E W\le2+L(4q^3-2q^4)\).

\item \textbf{Causal lower bound:} every deterministic local two-run strategy loses one exact-pair atom of mass \(h=q^2(1-q)^2\).  The product Bellman supermartingale, rather than a block-order assumption, proves \(\Pr[K\le2]\le(1-h)^L\) under arbitrary interleaving.

\item \textbf{Metric:} \(d(x,y)=\varepsilon\) for \(x<y\), \(d(x,y)=1\) otherwise, and depot legs \(1/2\) form a positive directed metric.  A service order with \(N\) active clients and \(K\) increasing runs has exact cost \(\varepsilon N+(1-\varepsilon)K\).

\item \textbf{Finite certificate:} \(q=1/100\), \(L=20000\), and \(\varepsilon=10^{-7}\) give \(3\operatorname{OPT}_{\rm adapt} -4\operatorname{OPT}_{\rm post}>0.24\).  The independent \(q=1/200,L=56000\) certificate separately gives \(3\operatorname{OPT}_{\rm adapt} -4\operatorname{OPT}_{\rm post}>0.10826\).

\end{itemize}

One wording qualification is important: the proof establishes \(\liminf \operatorname{OPT}_{\rm adapt}/\operatorname{OPT}_{\rm post} \ge3/2\).  Equality with \(3/2\) would require a matching adaptive upper bound, which is not needed for the strict construction.


\phantomsection
\section*{1. Shared diagnostics}
\addcontentsline{toc}{section}{1. Shared diagnostics}

\phantomsection
\subsection*{1.1 Additive replication — sound}
\addcontentsline{toc}{subsection}{1.1 Additive replication — sound}

\nolinkurl{composition-principles.md}, Section 1, is algebraically correct under its explicit hypothesis that both optimum values genuinely decompose as \(C+\sum_i P_i\) and \(C+\sum_i A_i\).  The note correctly warns that this hypothesis needs a ports/interleaving proof.


\phantomsection
\subsection*{1.2 Selector-first estimate — sound with one wording condition}
\addcontentsline{toc}{subsection}{1.2 Selector-first estimate — sound with one wording condition}

Section 3's estimate


\[
 D\sum_{v\in S}p_v
\]

is valid if \(D\) bounds every ordered distance between the depot and selector locations as well as between selectors.  In a fixed call order, movement occurs only along the active subsequence; there are \(\sum_v X_v\) such movements including the first depot-to-active movement. The note explicitly excludes the subsequent transition and depot return.


Suggested clarification: say “the directed diameter of \(S\cup\{r\}\)” every time this estimate is reused.  A diameter only within \(S\) does not bound the first active call.


\phantomsection
\subsection*{1.3 “Precisely” serving an active set — minor wording issue}
\addcontentsline{toc}{subsection}{1.3 “Precisely” serving an active set — minor wording issue}

Section 4 says that \(T_x\) must “visit precisely” the permanent and active clients.  In the shortest-path model, a generating walk may pass through inactive client locations.  The required property is:


\begin{ResearchQuote}

\(T_x\) explicitly serves every permanent and active client; its expanded paths may transit through any other vertex.

\end{ResearchQuote}

This does not affect the route-code diagnostic.


\phantomsection
\subsection*{1.4 Static clients do not open or close arcs — sound and important}
\addcontentsline{toc}{subsection}{1.4 Static clients do not open or close arcs — sound and important}

Section 6 is exactly consistent with the model.  Several speculative network templates would otherwise accidentally switch to a stochastic-edge model.


\phantomsection
\section*{2. Cayley digraphs}
\addcontentsline{toc}{section}{2. Cayley digraphs}

\phantomsection
\subsection*{2.1 Exact \(8/7\) word-order gadget — sound}
\addcontentsline{toc}{subsection}{2.1 Exact \(8/7\) word-order gadget — sound}

\nolinkurl{cayley-digraphs.md}, Section 1, survives a complete call-order audit.


\begin{itemize}

\item The listed generators realize the needed length-one and length-two marked distances.

\item Residual finiteness can simultaneously preserve marked-point distinctness and exclude the finite list of forbidden length-one equalities.

\item The posterior lower bounds count positive metric legs between distinct marked clients, not service obtained by transit.

\item Before the stochastic client \(s\) is called, the only observations are the deterministic outcomes of permanent clients, so the four placements of the \(s\)-call relative to \(b,h\) are exhaustive.

\item If a shortest expanded path passes through an uncalled marked vertex, it does not serve it.  The leg-count lower bounds therefore remain valid.

\end{itemize}

The order \(h,b,s\) is a legal policy: the calls to \(h,b\) force the first two movements, an inactive \(s\)-call leaves the policy at \(b\), and an active \(s\)-call forces the \(b\to s\) movement before the legal final return.


\phantomsection
\subsection*{2.2 Finite-group potential comment — harmless but should be sharpened}
\addcontentsline{toc}{subsection}{2.2 Finite-group potential comment — harmless but should be sharpened}

Section 2 says that a “homomorphism or weighted potential” \(\phi:\Gamma\to\mathbb R\) can certify word-metric lower bounds.  A group homomorphism from a finite group to \((\mathbb R,+)\) is necessarily trivial. A nonconstant *vertex potential* satisfying \(\phi(gs)-\phi(g)\le w_s\) can still be useful.


Suggested correction: remove “homomorphism” in the finite-group setting, or say that a homomorphism may instead target an ordered quotient before passing to a finite truncation.


No ratio or theorem depends on this sentence.


\phantomsection
\section*{3. Switching networks}
\addcontentsline{toc}{section}{3. Switching networks}

\phantomsection
\subsection*{3.1 Butterfly output-terminal theorem — sound}
\addcontentsline{toc}{subsection}{3.1 Butterfly output-terminal theorem — sound}

\nolinkurl{switching-networks.md}, Section 1, is correct.


\begin{itemize}

\item The only outgoing arc of an output is its return arc to \(r\).

\item Every route from \(r\) to an output uses one \(\alpha\)-arc and exactly \(k\) layered arcs.

\item Consequently every move between two distinct active output clients resets through \(r\) and costs \(C=\alpha+\beta+k\).

\item Inactive calls do not change the current output, so an execution with \(q\) active calls has realized cost exactly \(qC\), independent of the call order.

\end{itemize}

This proof does not rely on physically traversing a chosen butterfly route before an active call: the metric move forced by the active output call already has the claimed exact distance.


\phantomsection
\subsection*{3.2 Causal-prefix-tree observation — sound if tours mean client orders}
\addcontentsline{toc}{subsection}{3.2 Causal-prefix-tree observation — sound if tours mean client orders}

Section 3's gap-one observation is correct when each \(T_x\) is interpreted as a depot order of served clients and each realization-dependent branch calls the relevant selector.  It should not be interpreted as permission to voluntarily walk through nonclient switching vertices.  Metric movement between the successive called active clients automatically shortcuts those internal route vertices.


\phantomsection
\subsection*{3.3 General Beneš statement — correctly conditional}
\addcontentsline{toc}{subsection}{3.3 General Beneš statement — correctly conditional}

The extension in lines 80–83 is valid only for terminal constructions with a mandatory depot reset and equal depot-to-terminal cost.  The note states those conditions, so it is not a theorem about arbitrary Beneš uses.


\phantomsection
\section*{4. Projective-plane incidence graphs}
\addcontentsline{toc}{section}{4. Projective-plane incidence graphs}

\phantomsection
\subsection*{4.1 Potential identity — sound}
\addcontentsline{toc}{subsection}{4.1 Potential identity — sound}

\nolinkurl{projective-planes.md}, Section 2, correctly writes every directed incidence weight as a symmetric weight plus a vertex-potential difference.  The identity survives shortest-path minimization because every path between fixed endpoints has the same potential difference.  It also holds policy-by-policy: a realized adaptive call sequence begins and ends at \(r\), so all potential terms telescope.


\phantomsection
\subsection*{4.2 Singer Hamilton cycle — sound}
\addcontentsline{toc}{subsection}{4.2 Singer Hamilton cycle — sound}

Section 3's orbit construction gives a Hamilton cycle alternating through all points and lines.  The last translated line contains both \(p_{N-1}\) and \(p_0\), so the cyclic closure is valid.


\phantomsection
\subsection*{4.3 Exact posterior and adaptive values — sound}
\addcontentsline{toc}{subsection}{4.3 Exact posterior and adaptive values — sound}

Sections 4–6 correctly handle inactive lines as transit vertices.


The adaptive policy does not voluntarily traverse \(p_i\to\ell_i\to p_{i+1}\).  It calls \(\ell_i\), and:


\begin{itemize}

\item if active, the call moves to \(\ell_i\), after which the permanent call to \(p_{i+1}\) moves onward;

\item if inactive, it remains at \(p_i\), and the permanent call to \(p_{i+1}\) uses a metric distance no larger than the two-arc incidence path.

\end{itemize}

At the last line, the model permits the final return to the depot.  The expanded return may transit through the already-called \(p_0\).


The text twice says that an inactive point-to-point segment “costs” \(\alpha+\beta\); a shortcut could make it smaller.  Replace “costs” by “costs at most.”  The universal lower bound then forces equality for the total tour, so the theorem is unchanged.


\phantomsection
\subsection*{4.4 Expanded-walk lower bound — sound}
\addcontentsline{toc}{subsection}{4.4 Expanded-walk lower bound — sound}

A closed collection of incidence excursions based at \(p_0\) that contains all \(N\) permanent point endpoints has at least \(N\) point-to-point steps and therefore at least \(2N\) incidence arcs.  Passing through an uncalled point does not itself serve it, but only strengthens the need for an explicit point endpoint/service occurrence.  Multiple depot excursions do not help because the unique depot interface is \(p_0\).


\phantomsection
\section*{5. Generalized polygons}
\addcontentsline{toc}{section}{5. Generalized polygons}

\phantomsection
\subsection*{5.1 Potential obstruction — sound, but it is not by itself gap one}
\addcontentsline{toc}{subsection}{5.1 Potential obstruction — sound, but it is not by itself gap one}

\nolinkurl{generalized-polygons.md}, Section 3, correctly proves equivalence to the symmetric companion metric.  This removes the *directed* effect; it does not prove that an arbitrary non-Hamiltonian symmetric companion has clairvoyance gap one.  The note respects this distinction in Sections 3 and


\begin{enumerate}

\item The opening verdict should continue to keep the exact gap-one statement conditional on Hamiltonicity.

\end{enumerate}

\phantomsection
\subsection*{5.2 Tree posterior bound — sound}
\addcontentsline{toc}{subsection}{5.2 Tree posterior bound — sound}

Section 4's pruning leaves a tree containing every point.  Attaching each active outside line by one incidence edge creates a connected subgraph. Traversing every tree edge once in each direction is a valid posterior walk; inactive line vertices may be transit vertices.


\phantomsection
\subsection*{5.3 Universal alternating-walk lower bound — sound}
\addcontentsline{toc}{subsection}{5.3 Universal alternating-walk lower bound — sound}

Section 5 counts expanded incidence-arc occurrences.  A closed alternating walk has equally many point and line positions.  Each permanent point and active line must occur as a service endpoint somewhere in the expanded execution, so


\[
 H\ge2\max\{N_P,K\}.
\]

The potential telescoping then gives (5.1)–(5.2).  Interleaving and passages through uncalled vertices do not invalidate the count.


\phantomsection
\subsection*{5.4 Conditional Hamiltonian theorem — sound}
\addcontentsline{toc}{subsection}{5.4 Conditional Hamiltonian theorem — sound}

Section 6 uses a fixed sequence of actual client calls.  When a line is inactive, the next permanent point call performs the point-to-point movement; when active, the line call and point call split the same incidence segment. The final return handles the last line.  As in the projective-plane note, “costs” should be read as “costs at most” before applying the lower bound.


\phantomsection
\subsection*{5.5 Selector-first bound — sound}
\addcontentsline{toc}{subsection}{5.5 Selector-first bound — sound}

Section 7's policy is legal.  After all active line calls have been served, remove active lines from an optimal posterior client order and call the remaining permanent points in the resulting order.  If \(z\) is the last active selector and \(x\) the first permanent client, then


\[
 d(z,x)\le d(z,r)+d(r,x),
\]

so no voluntary repositioning is hidden in the proof.


The asymptotic statement after (7.4) needs the intended scale assumption made explicit: one needs \(\gamma=O(\alpha+\beta)\) and bounded ratios among \(\alpha,\beta,\gamma\), not merely an informal reference to “comparable” edge scales.


\phantomsection
\section*{6. High-girth lifts}
\addcontentsline{toc}{section}{6. High-girth lifts}

\phantomsection
\subsection*{6.1 Articulation-lift additivity — needs a policy-gluing paragraph}
\addcontentsline{toc}{subsection}{6.1 Articulation-lift additivity — needs a policy-gluing paragraph}

\nolinkurl{high-girth-lifts.md}, Section 2, has a sound lower-bound decomposition.  Once all metric moves are expanded, an excursion inside \(J_i-\{q_i\}\) starts and ends at the unique articulation \(q_i\).  Concatenating all such excursions gives a closed local service, while deleting them projects to a scaffold walk visiting every permanent \(q_i\).


The written adaptive upper bound is incomplete.  The matching policy cannot literally:


\begin{enumerate}

\item voluntarily finish a local tour by returning to \(q_i\);

\item walk along the scaffold to \(q_j\); and

\item start the next local service.

\end{enumerate}

Only active calls cause movements.  The repair is the Chapter 4 policy-gluing argument:


\begin{ResearchQuote}

Keep the local final path to \(q_i\), the next scaffold segment, and the next module's entry path as a pending generating walk.  Inactive calls leave it pending.  The next active client call shortcuts the whole pending walk in the metric.  Absorb the last pending suffix into the legal final depot return.

\end{ResearchQuote}

The same pending-walk argument realizes the scaffold walk despite revisits to already-called \(q_i\)'s.  Articulation also ensures that internal module walks cannot shorten scaffold distances.  With these sentences added, the exact formulas


\[
 C(S_N)+NP_J,\qquad C(S_N)+NA_J
\]

are justified.


The lower-bound simulation should likewise mention that outside activations are sampled as an independent seed and are revealed to the virtual ambient execution only at its virtual calls.  This is the same causality argument as Chapter 4's interrupted-services lemma.


\phantomsection
\subsection*{6.2 Local port-coupled constants — illustrative, not proved for the}
\addcontentsline{toc}{subsection}{6.2 Local port-coupled constants — illustrative, not proved for the}

ambient lift


Section 3, lines 113–133, specifies only a triangle and vague “permanent work behind” its ports.  The claims


\[
 P_{\rm local}=\tfrac32w,\qquad A_{\rm local}=2w
\]

are the Chapter 4 quotient constants only if:


\begin{itemize}

\item both port blocks are compulsory;

\item entry and exit are given the same open-service relaxation as Chapter 4;

\item outer lift arcs do not introduce a shorter trace; and

\item extra pieces are charged at the parent boundary.

\end{itemize}

As written, these are not values of the displayed ambient lift cell.  The surrounding text treats them as an uncoupled heuristic, but the wording “locally” and “give exactly” should be replaced by an explicit conditional.


\phantomsection
\subsection*{6.3 Ball packing statement — needs regularity}
\addcontentsline{toc}{subsection}{6.3 Ball packing statement — needs regularity}

The assertion that radius-\(R\) balls have \(\Theta((d-1)^R)\) vertices is valid for a \(d\)-regular lift below half the girth.  For an arbitrary bounded-degree lift it is only an upper bound without a minimum-degree assumption.  This does not affect the negative verdict.


\phantomsection
\section*{7. Ramanujan bigraphs}
\addcontentsline{toc}{section}{7. Ramanujan bigraphs}

\phantomsection
\subsection*{7.1 Euler fallback tour and adaptive pending sweep — sound}
\addcontentsline{toc}{subsection}{7.1 Euler fallback tour and adaptive pending sweep — sound}

\nolinkurl{ramanujan-bigraphs.md}, Section 2, gives a valid realization-wise walk of length


\[
 2|E|w+|A|w
\]

up to depot attachment.  It explicitly includes each active \(m_e\).


The adaptive implementation in lines 66–75 correctly uses a pending-walk argument.  Permanent vertices are called at their first marked occurrences; each midpoint is called at its marked reverse-edge occurrence.  The current position need not literally be the marked tail \(v\): the metric move to the next active client shortcuts the entire pending Euler segment.  Passing through an uncalled client on that segment does not serve it, but that client is still called at its own marked position later.


Suggested clarification: call this an upper bound for a fixed adaptive order, not a statement that the policy physically executes the Euler walk.


\phantomsection
\subsection*{7.2 Divergence correction — sound only for the retained-trace ledger}
\addcontentsline{toc}{subsection}{7.2 Divergence correction — sound only for the retained-trace ledger}

Lines 82–105 correctly calculate


\[
 b(x)=\pm(d-2A_x)
\]

for the directed multiset containing one locally preferred contracted trace for \textbf{every} edge.  Any Eulerian augmentation which retains that entire multiset must add at least \(\frac12\sum_x|b(x)|\) units of directed correction, and its expectation is \(\Omega(n\sqrt d)\) at \(p=1/2\).


This is not a lower bound on \(\operatorname{OPT}_{\rm post}\).  An optimum tour may omit inactive edge traces, use a different orientation, or route through other cells.  Section 2 recognizes the inactive-edge issue, but Section 4's phrase “the balance correction increases the posterior denominator” is too broad.


Suggested correction:


\begin{ResearchQuote}

The balance correction increases the cost of the all-edges independently-preferred candidate tour; it does not lower-bound the true posterior optimum.

\end{ResearchQuote}

\phantomsection
\subsection*{7.3 Separated-edge accounting — explicitly conditional and therefore sound}
\addcontentsline{toc}{subsection}{7.3 Separated-edge accounting — explicitly conditional and therefore sound}

Section 3 conditions the \(\tfrac32w,2w\) contributions on port boundaries strong enough to prove exact additivity.  Under that hypothesis the composition inequality gives at most \(4/3\).  It should not be cited later as an analyzed construction until the duplicated-port backbone and its ports/pieces proof are written.


\phantomsection
\subsection*{7.4 “At most one matched edge per vertex” — heuristic, not a universal}
\addcontentsline{toc}{subsection}{7.4 “At most one matched edge per vertex” — heuristic, not a universal}

charging theorem


Lines 143–149 explain why one perfect matching gives only \(n\) disjoint edge conflicts.  They do not prove that no more sophisticated bounded- congestion charging over several matchings exists.  The note correctly labels the desired summation as unjustified, so no theorem needs retraction.


\phantomsection
\section*{8. Algebraic lifts}
\addcontentsline{toc}{section}{8. Algebraic lifts}

\phantomsection
\subsection*{8.1 Product-choice permutation propositions — sound}
\addcontentsline{toc}{subsection}{8.1 Product-choice permutation propositions — sound}

\nolinkurl{algebraic-lifts.md}, Section 2, is correct as a statement about pointwise choices between two permutations.


After relabeling outputs, an output \(h\) has indegree


\[
 (1-X_h)+X_{\pi^{-1}(h)}.
\]

Every nonfixed point is defective with probability \(2p(1-p)\), giving


\[
 \mathbb E D=2p(1-p)|\operatorname{supp}\pi|.
\]

On each permutation cycle, holes and collisions alternate in equal number. Changing one input bit affects at most two defect indicators, so the stated linear high-probability conclusion follows from bounded differences when \(p\) is fixed and the support is linear.


These are combinatorial propositions.  They do not by themselves give a stochastic-TSP cost.


\phantomsection
\subsection*{8.2 Linear posterior repair claim — overstated}
\addcontentsline{toc}{subsection}{8.2 Linear posterior repair claim — overstated}

Section 3, especially lines 111–147, requires three assumptions not present in the displayed voltage layer:


\begin{enumerate}

\item output \(O_h\) must be identified with a state/input vertex having one compulsory outgoing preferred trace, or an equivalent balance condition must be imposed;

\item every “hole” port must be compulsory for the tour, rather than merely an unused output of the chosen pointwise map; and

\item the optimum posterior tour must be forced to retain every locally preferred trace.

\end{enumerate}

With distinct \(I_g,O_g\), a missing image of the pointwise map is not automatically a divergence defect of a circulation.  Even after identifying the layers, a stochastic-TSP optimum can omit an inactive trace unless permanent port work and a boundary lower bound force it.


The conditional statement that a circulation *containing all preferred traces* needs at least \(D/2\) additional unit entries is plausible once the identified-port model is made explicit.  It is not a lower bound on \(\operatorname{OPT}_{\rm post}\) for the template currently defined.


Suggested correction: retitle Section 3 “Conditional defect cost for a retained-trace circulation,” define the identified port graph, and avoid the word “posterior” until a tour-forcing lemma is supplied.


\phantomsection
\subsection*{8.3 The \(8/7\) ratio is not a rigorous obstruction to the gap}
\addcontentsline{toc}{subsection}{8.3 The \(8/7\) ratio is not a rigorous obstruction to the gap}

Lines 130–147 divide an assumed adaptive leading term \(2|\Gamma|w\) by the conditional retained-trace cost \(\tfrac74|\Gamma|w\).  This does not bound the actual gap in either direction:


\begin{itemize}

\item no adaptive upper bound of \(2|\Gamma|w\) is proved;

\item an adaptive policy may pay the same repair costs;

\item the posterior optimum may avoid the retained preferred traces; and

\item common scaffold/depot terms are absent.

\end{itemize}

The rigorous conclusion is only:


\begin{ResearchQuote}

the naïve posterior upper bound \(\tfrac32|\Gamma|w\) cannot be obtained by blindly concatenating all independently preferred permutation traces in the identified-port ledger.

\end{ResearchQuote}

The displayed \(8/7\) should be labeled a non-rigorous bookkeeping ratio or removed.  The Outcome and Verdict should not say that “the posterior pays” a linear repair term without the retained-trace qualification.


\phantomsection
\subsection*{8.4 Cheap correction-tour paragraph — incomplete construction}
\addcontentsline{toc}{subsection}{8.4 Cheap correction-tour paragraph — incomplete construction}

Lines 149–157 assert that cheap all-to-all connectors yield a tour of \(O((D+|\Gamma|)w)\).  The connector endpoints, compulsory ports, component joining rule, and depot attachment are not defined.  A fixed rail can provide a coarse walk, but its adaptive implementation again needs a fixed client order plus a pending-walk argument.  This paragraph should be treated as a template suggestion, not a proved a-posteriori bound.


\phantomsection
\subsection*{8.5 One global selector can use policy memory}
\addcontentsline{toc}{subsection}{8.5 One global selector can use policy memory}

Lines 169–177 say that one global client \(m_j\) “cannot remember” the incoming sheet.  The graph vertex cannot remember it, but the policy can.  With fan-in arcs \(I_g\to m_j\) and fan-out arcs \(m_j\to O_{ga_j}\), a posterior tour or adaptive policy can choose the appropriate outgoing destination using its known history.


The real obstruction is still serious: the same fan makes \(I_g\to m_j\to O_h\) available for many \(g,h\), collapsing sheet distances after shortest-path closure.  Suggested correction:


\begin{ResearchQuote}

A global selector is legal, but any fan-in/fan-out implementation must prove that policy memory cannot exploit the fan as a cheap arbitrary sheet switch.

\end{ResearchQuote}

The correlated-fiber warning in the next sentence is correct.


\phantomsection
\section*{9. Linear codes, designs, and locally testable structures}
\addcontentsline{toc}{section}{9. Linear codes, designs, and locally testable structures}

\phantomsection
\subsection*{9.1 Independent codeword proposition — sound}
\addcontentsline{toc}{subsection}{9.1 Independent codeword proposition — sound}

\nolinkurl{codes-designs-ltc.md}, Proposition 1, correctly shows that mutually independent nonconstant coordinates of a uniform linear-subspace word have full-cube support.  If all \(n\) coordinates are nonconstant and independent, the subspace is \(\mathbb F_2^n\).


The probability \(2^{k-n}\) in the independent-coordinate alternative is correct when \(k=\dim C\).


\phantomsection
\subsection*{9.2 Regular-design conditional theorem — sound}
\addcontentsline{toc}{subsection}{9.2 Regular-design conditional theorem — sound}

Theorem 2 is valid under its strong assumptions (P), (C), and (G). Conditioned on a complete causal history at the cut, every uncalled fair selector remains independent and fair.  If \(B_j\not\subseteq R\), its parity is therefore unpredictable with error \(1/2\).  The regularity count


\[
 tF(R)\le\rho|R|
\]

and the endpoint minimization in \(z=\mathbb E|R|\) are correct.


At line 188, write \(\mathbb E[M\mid\text{complete cut history}]\), not merely \(\mathbb E[M\mid R]\).  The displayed lower bound depends only on \(R\), so averaging yields the same theorem.  This is a precision correction, not a gap.


The theorem does not cover progressive/interleaved commitments; the note explicitly says so.


\phantomsection
\subsection*{9.3 Tanner potential theorem — sound with depot-potential clarification}
\addcontentsline{toc}{subsection}{9.3 Tanner potential theorem — sound with depot-potential clarification}

Section 4 correctly identifies uniform variable/check asymmetry as a vertex potential.  If the depot is symmetrically attached to one fixed Tanner vertex, set \(\phi(r)\) equal to that vertex's potential; the depot arcs then also have the required symmetric-plus-potential form.  This small definition should be stated explicitly when the fixed vertex is a check.


The theorem correctly claims equality with the symmetric companion values, not gap one.


\phantomsection
\subsection*{9.4 Bellman inequality has a mathematical typo and omits the boundary}
\addcontentsline{toc}{subsection}{9.4 Bellman inequality has a mathematical typo and omits the boundary}

Equation (5.3) is missing a plus sign.  It should read


\[
 \Phi(u,S)\le
 (1-p_v)\Phi(u,S\setminus\{v\})
 +p_v\bigl(d(u,v)+\Phi(v,S\setminus\{v\})\bigr)
\]

for every possible next call \(v\).


A full Bellman subsolution also needs the terminal condition


\[
 \Phi(u,\varnothing)\le d(u,r).
\]

Without the plus sign the displayed formula is not an inequality in the model's dynamic program.  The surrounding discussion is speculative, so no proved theorem is affected.


\phantomsection
\subsection*{9.5 Cover-menu all-active argument — sound under the fair-bit context}
\addcontentsline{toc}{subsection}{9.5 Cover-menu all-active argument — sound under the fair-bit context}

Section 6.2 uses that the all-active realization has positive probability. That is true under the Section 2 assumption \(p_i=1/2\).  For arbitrary probabilities it requires every \(p_i>0\).  The fixed universal route can be implemented by a fixed client call order; inactive calls disappear and metric moves shortcut the corresponding route segments.


\phantomsection
\subsection*{9.6 Cosmetic notation collision}
\addcontentsline{toc}{subsection}{9.6 Cosmetic notation collision}

The next lemma says “calling \(r\) systematic selectors,” while \(r\) elsewhere denotes the depot.  Rename this count \(h\) or \(q\).


\phantomsection
\section*{10. Cross-note conclusions}
\addcontentsline{toc}{section}{10. Cross-note conclusions}

\phantomsection
\subsection*{Results safe to cite as complete}
\addcontentsline{toc}{subsection}{Results safe to cite as complete}

\begin{itemize}

\item Cayley word-order gadget: gap exactly \(8/7\).

\item Butterfly output clients: gap exactly \(1\).

\item Uniform projective-plane incidence instance: gap exactly \(1\).

\item Uniform generalized-polygon incidence metric: equal to its symmetric companion; gap exactly \(1\) on a balanced Hamiltonian subclass.

\item Uniform Tanner orientation: equal to its symmetric companion.

\item Independent coordinates of a nontrivial proper linear codeword cannot be the activation law.

\item Pointwise choices between two distinct permutations cannot form a permutation for every independent bit vector.

\item Articulation lifts are additive after inserting the pending-walk gluing argument.

\end{itemize}

\phantomsection
\subsection*{Claims that should not be cited as stochastic-TSP bounds yet}
\addcontentsline{toc}{subsection}{Claims that should not be cited as stochastic-TSP bounds yet}

\begin{itemize}

\item the algebraic-lift \(\tfrac74\) posterior term and \(8/7\) ratio;

\item the coarse algebraic all-to-all repair tour;

\item exact \(\tfrac32w,2w\) values for an ambient port-coupled high-girth cell;

\item any posterior lower bound inferred from Ramanujan preferred-orientation imbalance; and

\item any multi-matching adaptive charge without a bounded-congestion ports/pieces proof.

\end{itemize}

\phantomsection
\subsection*{Recommended proof hygiene for the comparison note}
\addcontentsline{toc}{subsection}{Recommended proof hygiene for the comparison note}

Every candidate should be assigned one of four labels:


\begin{enumerate}

\item \textbf{exact model theorem;}

\item \textbf{conditional theorem with all hypotheses displayed;}

\item \textbf{candidate routing ledger, not an optimum bound;}

\item \textbf{speculative next lemma.}

\end{enumerate}

In particular, ratios formed from a candidate posterior *lower* cost and a hoped-for adaptive contribution should never be presented as a bound on \(\operatorname{OPT}_{\rm adapt}/\operatorname{OPT}_{\rm post}\).  A lower-gap construction needs an adaptive lower bound divided by a posterior upper bound; an upper obstruction needs an adaptive upper bound divided by a posterior lower bound.


\phantomsection
\section*{11. Finite buildings}
\addcontentsline{toc}{section}{11. Finite buildings}

\phantomsection
\subsection*{11.1 Potential and selector-first lemmas — sound}
\addcontentsline{toc}{subsection}{11.1 Potential and selector-first lemmas — sound}

\nolinkurl{finite-buildings.md}, Sections 2.1–2.2, correctly handles shortest-path closure and adaptive legality.  Potential differences telescope on every realized closed call sequence.  In the selector-first policy, movement occurs only along the active selector subsequence; after deleting those already served selectors from an offline client order, the first permanent call legally absorbs the transition from the last active selector.  When \(K=0\), the proof should simply start the permanent order at the depot; the harmless \(+\Delta\) slack in (2.5) still covers this case.


\phantomsection
\subsection*{11.2 Radial spherical template — bounds are sound}
\addcontentsline{toc}{subsection}{11.2 Radial spherical template — bounds are sound}

Sections 3.1–3.5 correctly identify radial gallery weights as a potential. The tree traversal in (3.2), expanded-walk count in (3.4), and fixed client order in (3.6) respect the model:


\begin{itemize}

\item a depth-first walk may transit inactive chambers, but explicitly includes every active chamber;

\item an expanded closed walk serving \(c_0\) and \(K\) distinct active chambers has at least \(K+1\) gallery arcs when \(K>0\); and

\item the fixed-order policy does not voluntarily follow gallery edges: the next active chamber call shortcuts the pending gallery segment, and the legal final return supplies the last segment.

\end{itemize}

The last point is implicit rather than stated, but the metric-distance proof of (3.6) is already sufficient.


\phantomsection
\subsection*{11.3 Apartment cardinality and one-hot calculation — one independence}
\addcontentsline{toc}{subsection}{11.3 Apartment cardinality and one-hot calculation — one independence}

qualification


Section 4.1's cardinality obstruction and Chernoff bound are correct.  Section 4.3 also correctly maximizes


\[
t p(1-p)^{t-1}\le\tfrac12.
\]

The product bound \(2^{-k}\), however, requires the \(k\) panels to use pairwise disjoint \textbf{selector-client sets}, not merely geometrically distinct panels.  Chambers incident to different panels can overlap, so “disjoint branch panels” should be replaced by “branch panels whose alternative selector chambers are disjoint,” or independence should be proved separately.


The sentence that one Bernoulli bit “chooses” a continuation should also be read only as a route label.  Activation does not enable or disable graph arcs; an explicit service gadget must make the two outcomes favor different client orders.


\phantomsection
\subsection*{11.4 Retraction lower bound — valid idea, but normalize units and depot}
\addcontentsline{toc}{subsection}{11.4 Retraction lower bound — valid idea, but normalize units and depot}

Section 5's nonexpanding projection is sound for the symmetric gallery length after the potential telescopes.  Equation (5.2) needs a convention:


\begin{itemize}

\item if \(\operatorname{TSP}_A(Q)\) uses unit gallery edges, the last term should be \(s\,\operatorname{TSP}_A(Q)\);

\item if it already uses edge length \(s\), the displayed formula is correct.

\end{itemize}

The apartment TSP must also be based at the projected depot chamber \(\rho(c_0)\).  These are normalization omissions, not failures of the retraction diagnostic.  The note correctly warns that transit/service multiplicity inside a retraction fiber is lost.


\phantomsection
\subsection*{11.5 Permanent-apartment ratio-one conclusion — overstated unless \(P\)}
\addcontentsline{toc}{subsection}{11.5 Permanent-apartment ratio-one conclusion — overstated unless \(P\)}

is a lower scale


Equation (6.1) is valid:


\[
\operatorname{OPT}_{\rm adapt}
\le\operatorname{OPT}_{\rm post}+\Delta(\mu_S+1).
\]

Lines 450–459 then say that \(\mu_S=o(P)\) forces ratio \(1+o(1)\), where the preceding paragraph describes a putative posterior tour “of cost \(P\).”  A posterior \textbf{upper bound} \(P\) cannot be used in the denominator of this ratio.  The conclusion requires


\[
\Delta(\mu_S+1)=o(\operatorname{OPT}_{\rm post}),
\]

or a separately proved lower bound \(\operatorname{OPT}_{\rm post}=\Omega(P)\).  Suggested correction: replace “\(\mu_S=o(P)\)” by the displayed condition involving the actual posterior optimum, and then state any convenient sufficient lower-bound hypothesis.


\phantomsection
\subsection*{11.6 Affine-quotient generic bounds — sound with the same scale condition}
\addcontentsline{toc}{subsection}{11.6 Affine-quotient generic bounds — sound with the same scale condition}

Equations (7.1)–(7.3) are valid provided every chamber adjacency has both directed arcs with weights in \([w_{\min},w_{\max}]\), as the gallery template intends.  Passing through an uncalled chamber does not serve it, but cannot reduce the number of distinct service endpoints needed for (7.2).


The selector conclusion around (7.4) additionally needs bounded \(w_{\max}\), depot scale, and


\[
\operatorname{OPT}_{\rm post}=\Omega(P_N)
\quad\text{or directly}\quad
\left(\sum_{v\in S}p_v\right)D_Xw_{\max}
=o(\operatorname{OPT}_{\rm post}).
\]

If “post cost \(P_N\)” is only the cost of a proposed backbone tour, the little-\(o\) claim does not follow.


\phantomsection
\subsection*{11.7 Cocycle and hybrid sections — appropriately speculative}
\addcontentsline{toc}{subsection}{11.7 Cocycle and hybrid sections — appropriately speculative}

Sections 7.2–10 correctly distinguish a global potential from a quotient cocycle, keep graph labels deterministic, allow inactive chambers as transit, and refrain from claiming a causal lower bound.  The proposed fiber-charged retraction inequality (11.1) is explicitly a next lemma, not an established ports decomposition.


\phantomsection
\subsection*{Finite-buildings verdict}
\addcontentsline{toc}{subsection}{Finite-buildings verdict}

The complete potential no-go and generic spherical/quotient bounds are safe to cite.  The two corrections needed before synthesis are:


\begin{enumerate}

\item qualify the \(2^{-k}\) one-hot probability by disjoint selector sets; and

\item replace the claims “\(\mu_S=o(P)\) implies ratio \(1+o(1)\)” and its affine analogue by conditions measured against a proved lower bound on \(\operatorname{OPT}_{\rm post}\).

\end{enumerate}
\end{document}
